\documentclass[11pt]{article}
\usepackage[letterpaper,margin=0.72in,top=0.82in,bottom=0.72in,headheight=22pt]{geometry}
\usepackage[T1]{fontenc}
\usepackage{lmodern}
\usepackage{microtype}
\usepackage[table]{xcolor}
\usepackage{amsmath,amssymb,mathtools}
\usepackage{array,longtable,booktabs,tabularx}
\usepackage[inline]{enumitem}
\usepackage[most]{tcolorbox}
\usepackage{fancyhdr}
\usepackage{hyperref}

\definecolor{Primary}{HTML}{17324D}
\definecolor{Accent}{HTML}{2C7A7B}
\definecolor{CueTint}{HTML}{EAF1F5}
\definecolor{NoteTint}{HTML}{F8FAFB}
\definecolor{Rule}{HTML}{B8C6CF}
\hypersetup{hidelinks,pdftitle={Chapter 15: Sequence Databases and Their Uses - The Mother Lode},pdfsubject={Cornell notes}}
\setlist[itemize]{leftmargin=1.25em,itemsep=1pt,topsep=1.5pt,parsep=0pt}
\setlength{\parindent}{0pt}
\setlength{\parskip}{3pt}
\renewcommand{\arraystretch}{1.16}
\emergencystretch=2em

\pagestyle{fancy}
\fancyhf{}
\fancyhead[L]{\sffamily\small\color{Primary} Chapter 15}
\fancyhead[R]{\sffamily\small\color{Primary} Sequence database search}
\fancyfoot[C]{\sffamily\small\color{Primary}\thepage}
\renewcommand{\headrulewidth}{0.5pt}
\renewcommand{\headrule}{\hbox to\headwidth{\color{Accent}\leaders\hrule height \headrulewidth\hfill}}

\newcolumntype{C}{>{\raggedright\arraybackslash}p{0.275\textwidth}}
\newcolumntype{N}{>{\raggedright\arraybackslash}p{0.655\textwidth}}
\newcommand{\cnrow}[2]{%
  \cellcolor{CueTint}\textbf{\sffamily\color{Primary}#1} &
  \cellcolor{NoteTint}#2 \\[4pt]\arrayrulecolor{Rule}\hline}

\begin{document}

\begin{tcolorbox}[
  enhanced,breakable,colback=Primary!4,colframe=Primary,boxrule=0.9pt,
  arc=2mm,left=3mm,right=3mm,top=2.5mm,bottom=2.5mm]
{\sffamily\color{Accent}\bfseries CHAPTER 15}\\[-1pt]
{\sffamily\LARGE\bfseries\color{Primary} Sequence Databases and Their Uses - The Mother Lode}

\vspace{2pt}
\textbf{Purpose.} Combine fast seed filters, biologically calibrated scoring, and statistical interpretation to search large protein and nucleotide databases.

\textbf{Learning objectives}
\begin{itemize}
  \item Separate sensitivity, selectivity, speed, and significance.
  \item Trace the filter-and-refine stages of FASTA and BLAST.
  \item Explain the evolutionary basis of PAM and BLOSUM matrices.
  \item Interpret motif, block, and database-size effects.
\end{itemize}

\textbf{Prerequisites.} Local alignment, hashing, word matching, protein alphabets, probability, and multiple alignment.
\end{tcolorbox}

\begin{longtable}{@{}C!{\color{Accent}\vrule width 0.8pt}N@{}}
\rowcolor{Primary}
\color{white}\sffamily\bfseries Cue / question &
\color{white}\sffamily\bfseries Notes \\ \hline
\endfirsthead
\rowcolor{Primary}
\color{white}\sffamily\bfseries Cue / question &
\color{white}\sffamily\bfseries Notes (continued) \\ \hline
\endhead
\cnrow{\S15.1: Why can database search succeed?}{A query sequence can reveal function, family membership, or structural clues through statistically unusual local similarity to annotated entries. Success depends on database quality and model calibration as much as on raw alignment score.}
\cnrow{\S15.2: What does database scale change?}{Collections grow rapidly, contain redundant and unevenly annotated records, and evolve over time. Search systems therefore need indexing, parallelism, stable identifiers, versioning, and significance measures that account for the amount of sequence examined.}
\cnrow{\S15.3: What are the algorithmic tensions?}{A full optimal local alignment against every entry is sensitive but expensive. Fast search uses a cheap filter to find promising regions and a more accurate extension or alignment to score them. Thresholds trade missed weak relatives against excessive candidates.}
\cnrow{\S15.4: What must a real search report?}{Report query and subject coordinates, raw and normalized scores, database version and size, expected false hits, alignment composition, and filtering settings. Repeated or low-complexity segments can produce high scores unrelated to the intended biological relationship.}
\cnrow{\S15.5: How does FASTA filter?}{Identify shared short words, accumulate hits along diagonals, join or rescore promising regions, and apply a more precise local alignment to top candidates. Diagonal clustering exploits the fact that an ungapped local similarity produces several word hits with a consistent offset.}
\cnrow{What is the FASTA profile?}{\textbf{Input/output:} query, database, ranked local similarities. \textbf{Preprocess:} word lookup for the query. \textbf{Filter:} high-density diagonal hits. \textbf{Refine:} optimized rescoring or local DP. \textbf{Tradeoff:} shorter words improve sensitivity but increase candidate work.}
\cnrow{\S15.6: How does BLAST seed and extend?}{Generate query words whose substitution-matrix score exceeds a threshold, locate database word hits, and extend each hit in both directions while the running score remains promising. Gapped refinement may follow. The high-scoring segment and statistical model support ranked reporting.}
\cnrow{What is the BLAST profile?}{\textbf{Invariant:} every active extension is anchored by a qualifying word hit. \textbf{Controls:} word length, neighborhood threshold, drop-off, gap parameters, and composition filters. \textbf{Cost:} dominated by seed hits and extensions rather than every possible full alignment.}
\cnrow{\S15.7: What are PAM matrices?}{Estimate accepted amino-acid substitutions from closely related proteins, convert them into log-odds scores, and extrapolate a family of matrices for different evolutionary distances. Matrix choice changes which conservative substitutions receive positive evidence.}
\cnrow{\S15.8: What does PROSITE represent?}{Curated patterns and profiles describe characteristic family or functional sites. A motif hit is compact and interpretable but may be too strict for divergent members or too common when the motif is short. Context and additional evidence remain necessary.}
\cnrow{\S15.9: What are BLOCKS?}{Blocks are ungapped conserved regions assembled from related proteins. They summarize position-specific variation without requiring a full-length alignment and provide data for family-specific scoring and matrix construction.}
\cnrow{\S15.10: How are BLOSUM matrices built?}{Count substitutions in conserved alignment blocks and cluster sequences above a chosen identity threshold to reduce sampling bias. Lower-numbered matrices represent more aggressively clustered, typically more divergent comparisons; the score is a log-odds contrast to background frequencies.}
\cnrow{PAM versus BLOSUM?}{PAM begins with an evolutionary model from close sequences and extrapolates distance. BLOSUM counts replacements directly in conserved blocks after identity-based clustering. Their numbering conventions run in opposite practical directions, so matrix selection should be deliberate.}
\cnrow{\S15.11: What else affects interpretation?}{Database size changes expected hit counts; amino-acid composition, repeats, low complexity, transmembrane regions, and biased query families alter null behavior. Masking, composition adjustment, reciprocal searches, domain databases, and multiple evidence sources improve reliability.}
\cnrow{\S15.12: What should practice emphasize?}{Trace seed, diagonal, extension, and refinement stages; compare matrix choices; and interpret a hit using score, expected frequency, coverage, identity, conserved motifs, and alternative explanations rather than one number.}
\end{longtable}


\begin{tcolorbox}[enhanced,breakable,title={Chapter synthesis},
  colback=Accent!5,colframe=Accent,fonttitle=\small\sffamily\bfseries,fontupper=\footnotesize,
  boxsep=0.6mm,left=1.5mm,right=1.5mm,top=1mm,bottom=1mm,before skip=3pt,after skip=3pt]
Large-scale sequence search is a calibrated filter-and-refine problem. FASTA clusters exact word evidence on diagonals; BLAST extends high-scoring neighborhood seeds; PAM and BLOSUM supply log-odds substitution evidence; motifs and blocks add family-specific structure. A ranked hit is meaningful only with database, composition, coverage, and statistical context.
\end{tcolorbox}

\begin{tcolorbox}[enhanced,breakable,title={Key takeaways},
  colback=white,colframe=Primary,fonttitle=\small\sffamily\bfseries,fontupper=\footnotesize,
  boxsep=0.6mm,left=1.5mm,right=1.5mm,top=1mm,bottom=1mm,before skip=3pt,after skip=3pt]
\begin{itemize}
  \item Sensitivity and speed are controlled by filter thresholds.
  \item FASTA uses word hits concentrated on diagonals.
  \item BLAST uses qualifying seeds and score-controlled extension.
  \item Substitution matrices encode log-odds evidence, not physical distance directly.
  \item PAM extrapolates a model; BLOSUM counts conserved-block substitutions.
  \item Database size and sequence composition affect significance.
  \item Coverage and biological context matter alongside the score.
\end{itemize}
\end{tcolorbox}

\begin{tcolorbox}[enhanced,breakable,title={Algorithm selection / when to use this},
  colback=Primary!3,colframe=Primary,fonttitle=\small\sffamily\bfseries,fontupper=\footnotesize,
  boxsep=0.6mm,left=1.5mm,right=1.5mm,top=1mm,bottom=1mm,before skip=3pt,after skip=3pt]
\begin{itemize}
  \item Use a sensitive local alignment for a small candidate set.
  \item Use FASTA- or BLAST-style filters for large databases.
  \item Use motif databases for interpretable conserved-site evidence.
  \item Choose a substitution matrix suited to expected divergence.
  \item Use composition filtering when biased or low-complexity regions dominate hits.
\end{itemize}
\end{tcolorbox}

\begin{tcolorbox}[enhanced,breakable,title={Self-test questions},
  colback=white,colframe=Accent,fonttitle=\small\sffamily\bfseries,fontupper=\footnotesize,
  boxsep=0.6mm,left=1.5mm,right=1.5mm,top=1mm,bottom=1mm,before skip=3pt,after skip=3pt]
\begin{itemize}
  \item Why do diagonal word-hit clusters suggest an ungapped similarity?
  \item Which parameters trade BLAST sensitivity for speed?
  \item How do PAM and BLOSUM construction philosophies differ?
  \item Why does database growth change expected false-hit counts?
  \item What evidence should accompany a reported high-scoring match?
  \item Why can a short motif produce misleading hits?
\end{itemize}
\end{tcolorbox}

{\footnotesize
\textbf{Terms and notation to remember.} seed, neighborhood word, diagonal, extension, drop-off, HSP, log odds, PAM, PROSITE, BLOCKS, BLOSUM, low complexity, expected hit count.

\textbf{Connections.} Database-search filters instantiate the exclusion principles of approximate matching; profiles and conserved blocks come from multiple alignments.
}

\end{document}
